\newpage
\section{Wstęp}
Dynamiczny rozwój technologii bezprzewodowych otwiera nowe możliwości monitorowania obiektów w 
środowiskach pozbawionych tradycyjnej infrastruktury sieciowej. Niniejsza praca eksploruje 
potencjał wspomnianych technologii w kontekście praktycznego zastosowania - śledzenia zwierząt na terenach wiejskich.

Praca dotyczy dwóch powiązanych dziedzin teleinformatyki. Pierwszą jest bezprzewodowa komunikacja
radiowa dalekiego zasięgu w obszarach pozbawionych infrastruktury sieciowej - czyli tam, gdzie nie ma
zasięgu sieci komórkowej ani punktów dostępowych WiFi. Szczególnie interesującą klasą rozwiązań
w tym kontekście są sieci radiowe typu mesh, w których każdy węzeł pełni jednocześnie rolę nadajnika,
odbiornika i przekaźnika, co eliminuje konieczność istnienia zewnętrznej infrastruktury.

Drugą dziedziną jest tworzenie aplikacji mobilnych integrujących dane z zewnętrznych urządzeń -
w tym przypadku urządzeń radiowych - i prezentujących je użytkownikowi w sposób użyteczny
w warunkach terenowych. Szczególnym zagadnieniem jest tu geofencing, czyli wykrywanie przekroczenia
przez obiekt wirtualnej granicy geograficznej wyłącznie w oparciu o dane lokalizacyjne i obliczenia
wykonywane lokalnie na urządzeniu mobilnym.

\subsection*{Motywacja}

Bezpośrednią inspiracją dla niniejszej pracy była obserwacja zachowania psa myśliwskiego rasy 
Mały Gończy Gaskoński w warunkach terenowych, gdzie silny
instynkt łowczy psa, połączony z otwartą przestrzenią graniczącą z terenami bogatymi w dziką
zwierzynę, prowadził do regularnych ucieczek za zwierzyną. Budowa ogrodzenia nie wchodziła w grę
ze względów krajobrazowych i prawnych, a trening odwołania okazywał się niewystarczający, ponieważ
pies pędząc w pogoni przestaje reagować na komendy.

Tak postawiony problem wymaga rozwiązania technicznego: systemu, który umożliwi śledzenie położenia
psa w czasie rzeczywistym i alarmowanie właściciela, gdy zwierzę oddali się poza bezpieczny obszar.
Ograniczeniem jest również brak zasięgu sieci komórkowej na tamtejszych terenach.

\subsection*{Cel pracy}

\textbf{Celem pracy inżynierskiej jest zaprojektowanie i implementacja systemu monitorowania
położenia zwierząt w czasie rzeczywistym, działającego w obszarach pozbawionych pokrycia sieci
komórkowej, wraz z aplikacją mobilną umożliwiającą wizualizację położenia na mapie, definiowanie
wirtualnych stref bezpieczeństwa oraz powiadamianie użytkownika o ich przekroczeniu.}
